% Generated draft: requires mathematical and source review.
\documentclass[11pt]{article}
\usepackage[margin=25mm]{geometry}
\usepackage{fontspec}
\setmainfont{Latin Modern Roman}
\usepackage{amsmath,amssymb,mathtools,unicode-math}
\setmathfont{Latin Modern Math}
\usepackage{enumitem,needspace,xurl,hyperref}
\setlength{\parindent}{0pt}
\setlength{\parskip}{5pt}
\setlength{\emergencystretch}{3em}
\newcommand{\N}{\mathbb{N}}
\newcommand{\Z}{\mathbb{Z}}
\newcommand{\Q}{\mathbb{Q}}
\newcommand{\R}{\mathbb{R}}
\newcommand{\C}{\mathbb{C}}
\newcommand{\F}{\mathbb{F}}
\newcommand{\A}{\mathbb{A}}
\newcommand{\PP}{\mathbb{P}}
\newcommand{\E}{\mathbb{E}}
\newcommand{\eps}{\varepsilon}
\newcommand{\dd}{\,\mathrm{d}}
\begin{document}
\section*{UnsolvedMath ID: 30002557}
\textbf{Source problem code:} catalogue entry
\subsection{Definitions and mathematical statement}
For an integer \(g\geq 2\), let \(\mathcal M_g\) denote the moduli object of smooth, connected, compact Riemann surfaces of genus \(g\), with no marked points. The supplied record does not specify whether this means the coarse complex analytic moduli space or the complex analytic moduli orbifold.

For a complex analytic space or orbifold \(X\), one possible precise convention for Dolbeault cohomological dimension is
\[
\operatorname{cohdim}_{\mathrm{Dol}}(X)
:=\sup\bigl\{q\in\Z_{\geq0}: H^q(X,\mathcal F)\neq0
\text{ for some coherent analytic }\mathcal O_X\text{-module }\mathcal F\bigr\},
\]
where \(H^q\) denotes analytic sheaf cohomology. This convention is provisional: the supplied record does not define the invariant or identify its coefficient category, and its agreement with the cited formulation requires source verification.

\textbf{Conjecture attributed to Looijenga in the supplied record.} For every integer \(g\geq2\),
\[
\operatorname{cohdim}_{\mathrm{Dol}}(\mathcal M_g)=g-2.
\]
Under the preceding coherent-analytic-sheaf convention and a fixed interpretation of \(\mathcal M_g\), this means that, for every integer \(g\geq2\),
\[
H^q(\mathcal M_g,\mathcal F)=0
\quad\text{for every coherent analytic sheaf }\mathcal F
\text{ and every integer }q>g-2,
\]
and there exists a coherent analytic sheaf \(\mathcal F_g\) on \(\mathcal M_g\) such that
\[
H^{g-2}(\mathcal M_g,\mathcal F_g)\neq0.
\]
The boundary case \(g=2\) is included. These expanded quantifiers describe the stated provisional convention; the supplied evidence does not establish that this is the exact intended definition. The conjecture is retained as a source-attributed question, not asserted as a theorem.
\subsection{Short English statement}
The supplied conjecture asks whether the Dolbeault cohomological dimension of the moduli space of compact genus-\(g\) Riemann surfaces is exactly \(g-2\) for every integer \(g\geq2\). The precise moduli interpretation and definition of Dolbeault cohomological dimension require verification. The supplied dated curation labels the problem partially solved.
\subsection{Sources}
\begin{itemize}
\item Local UnsolvedMath registry, supplied record with permanent ID \(30002557\), title ``Looijenga's Dolbeault Dimension Conjecture,'' source problem code \texttt{catalogue entry}, in the Oberwolfach Reports Open Problems collection. The record attributes the conjecture to Looijenga and lists proposed year 2014; no separate proposer is supplied.
\item \emph{Flat Surfaces and Dynamics on Moduli Space} (2014), Oberwolfach Reports, report 2014/15, pp.~869--941, DOI \texttt{10.4171/owr/2014/15}. Supplied report provenance and literature link: \url{https://doi.org/10.4171/owr/2014/15}. No problem-specific page is supplied.
\item Mondello, \emph{International Journal of Mathematics}. Supplied literature link: \url{https://doi.org/10.1142/S0129167X16400085}. The record supplies no article title or theorem locator.
\item Fullarton--Putman, arXiv. Supplied literature link: \url{https://arxiv.org/abs/1610.03768}. The record supplies no article title or theorem locator.
\item Literature assessment dated exactly \texttt{2026-08-21}: the supplied curation states that the equality is known in low genus, that general upper and lower bounds are available, and that the equality for arbitrary genus has not been proved. Its status label is \texttt{partially\_solved}. This is supplied dated curation derived from the OWR record, not an independently verified literature assessment; the record explicitly requests independent verification of time-sensitive status and cited claims.
\item The registry labels the statement \texttt{corrected\_verified} and reports manual semantic, typographic, and LaTeX review of a meaning-preserving cleanup. This records supplied editorial review, not independent verification of the conjecture or its mathematical conventions.
\end{itemize}
\end{document}
